\MulaiModul{05}{Perulangan}

\subsection{Tujuan}

\begin{tujuan}
Setelah menyelesaikan modul ini, mahasiswa mampu:
\begin{enumerate}
    \item menjelaskan kerja \texttt{LOOP} dan peran \texttt{CX}
          sebagai pencacah;
    \item menulis perulangan dengan cacah tertentu memakai
          \texttt{LOOP};
    \item menulis perulangan yang setara memakai \texttt{DEC} dan
          \texttt{JNZ} tanpa \texttt{LOOP};
    \item menelusuri perubahan \texttt{CX}, \texttt{ZF}, dan karakter
          yang dicetak pada setiap iterasi;
    \item merakit, menjalankan, dan memeriksa perulangan dengan
          emu8086 atau DEBUG.
\end{enumerate}
\end{tujuan}

\subsection{Dasar singkat}

\subsubsection*{LOOP dan pencacah CX}

\texttt{LOOP} memakai \texttt{CX} sebagai pencacah. Instruksi ini
mengurangi \texttt{CX} dengan satu, lalu melompat ke label tujuan
bila \texttt{CX} belum nol. Bila \texttt{CX} sudah nol, eksekusi
berlanjut ke instruksi sesudah \texttt{LOOP}. Karena pengurangan
dilakukan lebih dahulu, badan perulangan selalu dijalankan paling
sedikit satu kali.

Sebelum label perulangan, \texttt{CX} harus diisi dengan cacah yang
diinginkan. Bila \texttt{CX} bernilai nol pada awal perulangan,
pengurangan akan menghasilkan \texttt{FFFFh} dan perulangan berjalan
65.536 kali. Karena itu, tetapkan \texttt{CX} lebih besar dari nol.
Badan perulangan juga tidak boleh mengubah \texttt{CX}, kecuali
memang ingin mengubah cacahnya.

\subsubsection*{DEC dan JNZ sebagai pembanding}

Perulangan yang sama dapat dibentuk dengan dua instruksi:
\texttt{DEC CX} lalu \texttt{JNZ ULANG}. \texttt{DEC} mengurangi
\texttt{CX} dengan satu dan menyalakan \texttt{ZF} ketika hasilnya
nol. \texttt{JNZ} melompat selama \texttt{ZF} bernilai nol. Hasil
akhirnya sama dengan \texttt{LOOP}, tetapi pemeriksaan flag terlihat
jelas karena ditulis sebagai dua instruksi. Perbandingan kedua bentuk
ditunjukkan pada Tabel~\ref{tab:m05-bentuk}.

\begin{table}[H]
    \centering
    \caption{Dua bentuk perulangan dengan pencacah \texttt{CX}.}
    \label{tab:m05-bentuk}
    \small
    \begin{tabular}{@{}lll@{}}
        \toprule
        \textbf{Bentuk} & \textbf{Langkah} & \textbf{Ciri} \\
        \midrule
        \texttt{LOOP ULANG} & \texttt{CX}$\leftarrow$\texttt{CX}$-1$;
        lompat bila \texttt{CX}$\neq0$ & ringkas, satu instruksi \\
        \texttt{DEC CX}; \texttt{JNZ ULANG} &
        kurangi, lalu periksa \texttt{ZF} & flag terlihat, dua instruksi \\
        \bottomrule
    \end{tabular}
\end{table}

\subsubsection*{Deretan karakter}

Perulangan dapat mencetak deretan karakter tanpa menelusuri array.
Karakter pertama diletakkan pada \texttt{DL}, dicetak dengan
\texttt{AH=02h}, lalu \texttt{INC DL} memajukannya ke karakter
berikutnya. Setiap iterasi mencetak satu karakter dan memakai
\texttt{LOOP} atau \texttt{DEC}/\texttt{JNZ} untuk kembali ke awal
badan perulangan. Jejak lima iterasi pertama ditunjukkan pada
Tabel~\ref{tab:m05-jejak}.

\begin{table}[H]
    \centering
    \caption{Jejak perulangan lima kali dengan \texttt{CX=5} dan
    \texttt{DL='A'}.}
    \label{tab:m05-jejak}
    \small
    \begin{tabular}{@{}cccc@{}}
        \toprule
        \textbf{Iterasi} & \textbf{CX sebelum} & \textbf{DL} &
        \textbf{CX sesudah} \\
        \midrule
        1 & 5 & A & 4 \\
        2 & 4 & B & 3 \\
        3 & 3 & C & 2 \\
        4 & 2 & D & 1 \\
        5 & 1 & E & 0, perulangan berhenti \\
        \bottomrule
    \end{tabular}
\end{table}

Pola \texttt{LOOP} sesuai dengan perulangan \texttt{for} yang
cacahnya diketahui. Perulangan \texttt{while} memeriksa syarat
sebelum badan dijalankan, sehingga badan dapat dilewati sama sekali;
pola itu tidak dipakai pada modul ini. Modul ini tidak memakai
\texttt{SI}, \texttt{[SI]}, \texttt{PROC}, \texttt{PUSH}, atau
penelusuran array.

\begin{rujukan}
Rujukan singkat: Gunawan, Bagian 11.3; praktikum \texttt{bab3.tex}
bagian LOOP; buku ajar, \texttt{chapters/13} bagian LOOP, while, dan
for.
\end{rujukan}

\subsection{Alat dan bahan}

\begin{itemize}
    \item Komputer dengan emu8086 terpasang.
    \item DOSBox dengan TASM, TLINK, dan DEBUG.
    \item Berkas \texttt{code/m05-loop-karakter.asm} dan
          \texttt{code/m05-dec-jnz.asm}.
    \item Lembar jurnal praktikum.
\end{itemize}

\subsection{Contoh program}

Program pertama mencetak lima karakter dengan \texttt{LOOP}.
\texttt{CX} diisi 5, \texttt{DL} diisi karakter pertama, lalu
\texttt{INC DL} memajukan karakter pada setiap iterasi.

\lstinputlisting[
    caption={Deretan karakter dengan \texttt{LOOP}.},
    label={lst:m05-loop}
]{code/m05-loop-karakter.asm}

Program kedua mencetak deretan yang sama tanpa \texttt{LOOP}.
Perulangan diakhiri dengan \texttt{DEC CX} dan \texttt{JNZ ULANG}.

\lstinputlisting[
    caption={Deretan karakter yang sama dengan \texttt{DEC} dan
    \texttt{JNZ}.},
    label={lst:m05-dec-jnz}
]{code/m05-dec-jnz.asm}

\subsection{Langkah kerja}

\begin{enumerate}
    \item Siapkan folder kerja, lalu salin kedua berkas \texttt{.asm}
          ke dalamnya.
    \item Rakit dan tautkan program pertama. Jalankan kedua perintah
          berikut pada baris terpisah:
          \begin{perintah}
tasm m05-loop-karakter.asm\\
tlink m05-loop-karakter.obj
\end{perintah}
          Jalankan hasilnya dan catat deretan karakter yang tercetak.
    \item Ubah \texttt{MOV CX,5} menjadi \texttt{MOV CX,8}, rakit
          ulang, lalu jalankan. Catat perubahan panjang deretan.
    \item Jalankan program pertama dengan \emph{single step} pada
          emu8086. Catat nilai \texttt{CX} dan \texttt{DL} pada
          setiap iterasi, lalu bandingkan dengan
          Tabel~\ref{tab:m05-jejak}.
    \item Rakit dan tautkan program kedua dengan perintah yang sama,
          lalu jalankan. Pastikan deretannya sama dengan program
          pertama.
    \item Jalankan program kedua dengan \emph{single step}. Amati
          nilai \texttt{ZF} sesudah \texttt{DEC CX} dan keputusan
          \texttt{JNZ}. Jelaskan mengapa hasilnya setara dengan
          \texttt{LOOP}.
\end{enumerate}

\subsection{Tugas}

\begin{enumerate}
    \item Jelaskan urutan kerja \texttt{LOOP}: kapan \texttt{CX}
          dikurangi dan kapan lompatan dilakukan.
    \item Jelaskan mengapa \texttt{DEC CX} diikuti \texttt{JNZ} memberi
          hasil yang sama dengan \texttt{LOOP}.
    \item Jelaskan apa yang terjadi bila \texttt{CX} bernilai nol
          sebelum perulangan dimulai, dan mengapa cacah harus lebih
          besar dari nol.
    \item Tentukan keluaran program bila \texttt{CX=3} dan
          \texttt{DL} diisi \texttt{'X'}.
    \item Bandingkan bentuk \texttt{LOOP} dan \texttt{DEC}/\texttt{JNZ}
          dari segi panjang instruksi dan kejelasan pemeriksaan flag.
\end{enumerate}

\subsection{Hasil yang dikumpulkan}

\begin{enumerate}
    \item Kedua berkas \texttt{.asm} beserta versi modifikasinya.
    \item Tangkapan layar keluaran kedua program.
    \item Tabel jejak iterasi berisi \texttt{CX}, \texttt{DL}, dan
          karakter yang tercetak untuk kedua program.
    \item Lembar jurnal berisi jawaban langkah kerja dan tugas.
\end{enumerate}
